% Generated draft: requires mathematical and source review.
\documentclass[11pt]{article}
\usepackage[margin=25mm]{geometry}
\usepackage{fontspec}
\setmainfont{Latin Modern Roman}
\usepackage{amsmath,amssymb,mathtools,unicode-math}
\setmathfont{Latin Modern Math}
\usepackage{enumitem,needspace,xurl,hyperref}
\setlength{\parindent}{0pt}
\setlength{\parskip}{5pt}
\setlength{\emergencystretch}{3em}
\newcommand{\N}{\mathbb{N}}
\newcommand{\Z}{\mathbb{Z}}
\newcommand{\Q}{\mathbb{Q}}
\newcommand{\R}{\mathbb{R}}
\newcommand{\C}{\mathbb{C}}
\newcommand{\F}{\mathbb{F}}
\newcommand{\A}{\mathbb{A}}
\newcommand{\PP}{\mathbb{P}}
\newcommand{\E}{\mathbb{E}}
\newcommand{\eps}{\varepsilon}
\newcommand{\dd}{\,\mathrm{d}}
\begin{document}
\section*{UnsolvedMath ID: 20003323}
\textbf{Source problem code:} catalogue entry
\subsection{Definitions and mathematical statement}
This item is commentary, not an independently stated conjecture. Its exact content says that Conjecture 36 is based on computations in [10, 11], that the supremum $m_2(Q)$ need not be isolated (see [71]), and then opens Section 5.7 with a question communicated by D. Kleinbock. Neither Conjecture 36, the definition of $m_2(Q)$, nor the following Question 37 is contained in this record. The Kleinbock attribution belongs to that following question, not automatically to the preceding non-isolation claim.

The supplied research aid interprets the first part as commentary on a second inhomogeneous minimum and reports that Godwin's cited 1963 result evaluates the second minimum for $Q(x,y)=x^2-23y^2$ and proves it non-isolated. The exact definition, numerical value, and complete citation are not supplied here, so no value is asserted.

A separate abstract formulation of the aid's topological reduction is as follows. Let $X$ be a compact metric space and $f:X\to\R$ upper semicontinuous. Set
\[
 M=\max_X f,\qquad A=\{x:f(x)=M\},\qquad
 m=\sup_{x\in X\setminus A} f(x),
\]
and assume $X\setminus A\ne\varnothing$ and $m<M$; the latter is the isolated-top-value hypothesis. Then the second level $m$ is attained on $X\setminus A$ if and only if there exist $\varepsilon>0$ and a sequence $x_j\in X\setminus A$ with
\[
 f(x_j)\longrightarrow m,\qquad
 \operatorname{dist}(x_j,A)\geq\varepsilon\quad\hbox{for every }j.
\]
In particular, if $m$ is not attained, every such maximizing sequence satisfies $\operatorname{dist}(x_j,A)\to0$. If a group $H$ acts on $X$, $f$ is $H$-invariant, and $m$ is non-isolated \emph{from below} in the value set, meaning
\[
 \forall\varepsilon>0\ \exists x\in X:
             m-\varepsilon<f(x)<m,
\]
then one may choose points with distinct values strictly increasing to $m$; these points belong to pairwise distinct $H$-orbits. This companion statement does not identify $X,f,H$ with the undefined quadratic-form objects in the source fragment. Attainment and isolation are separate properties. The registry marks this record open, while the aid classifies the extracted text as context-only/not-a-problem.
\subsection{Short English statement}
The entry comments on computations and the possible non-isolation of a second inhomogeneous minimum, then introduces a different question. Its supplied companion lemma explains when a second level of an upper-semicontinuous function is attained and why distinct nearby invariant values give distinct group orbits. The actual surrounding conjecture and the definition of $m_2(Q)$ are missing.
\subsection{Sources}
\begin{itemize}
\item Local UnsolvedMath registry, permanent ID 20003323, exact code catalogue entry; \texttt{aim-other-notes.json}, notes[129], commentary after Conjecture 36 and transition to Section 5.7/Question 37.
\item American Institute of Mathematics, \emph{Emerging applications of measure rigidity}, \url{https://aimath.org/WWN/measrigid/measrigid.pdf}. Preserve the source's numerical citations [10, 11] for computations and [71] for non-isolation; bibliography entries are not provided. D. Kleinbock is credited for the following question.
\item Supplied research attempt 2: context\_only, problem status not\_a\_problem, rigor label proved, novelty confidence low. It reports Godwin's 1963 result for $x^2-23y^2$ and the compact-space lemma. This machine-generated aid is not peer reviewed and requires independent verification before citation; a complete reference to Godwin is absent.
\end{itemize}
\end{document}
